\section{Related Work}
\label{sec:related-work}

\paragraph{Long-horizon RL benchmarks.}
The Arcade Learning Environment~\citep{bellemare2013arcade} provided
the dominant benchmark for value-based and policy-gradient deep RL
methods through the late 2010s, but its tasks rarely exceed a few
thousand frames per episode and have dense, well-shaped reward
signals.  Procgen~\citep{cobbe2020procgen} added procedural generation
to stress generalisation but kept per-episode horizons short.  More
recent open-ended benchmarks---Craftax~\citep{matthews2024craftax} and
MiniHack~\citep{samvelyan2021minihack}---move toward the
sparse-reward, long-horizon regime, but do so in synthetic worlds whose
dynamics are simple enough to render in compute-light JAX.  The
present work targets the same regime via a real, hand-designed game
environment, where the causal chains and reward sparsity are products
of human design rather than procedural generation.

\paragraph{Pok\'emon Red as an RL test bed.}
The Pok\'emon-RL community has produced influential code releases and
blog-post-quality empirical results on Pok\'emon
Red~\citep{whitney2023pokemonrl,pokemonred_experiments}, demonstrating
that PPO with a pixel-based observation can learn meaningful behaviour
within the first few towns of the game.  These projects have been
crucial for showing the environment is tractable, but they do not
present controlled comparisons across observation treatments, do not
report bootstrap-based aggregate statistics, and have not been
released as installable Gymnasium environments.  We build directly on
their foundation while adding the methodological scaffolding required
for publication-ready empirical claims.

\paragraph{Observation representations in RL.}
The pixel-vs-state question has been studied in narrower contexts:
\citet{anand2019unsupervised} showed that learned state representations
on Atari preserve much of the underlying game state, suggesting that
pixel agents \emph{re-derive} information that a symbolic agent could
read directly.  \citet{kornblith2019cka} introduced CKA as a tool for
comparing internal representations across networks, which we plan to
apply in follow-up work to ask whether pixel and symbolic agents
converge to similar internal features.  To our knowledge, no prior
work performs a matched-architecture, matched-reward, matched-compute
comparison of pixel, symbolic, and hybrid observation streams in a
single long-horizon RL task.

\paragraph{Recurrent policies under partial observability.}
Pok\'emon Red is partially observable from any single frame: the
overworld map, the player's party, and the event-flag state of the
world are not visible simultaneously.  Recurrent policies are the
canonical response to partial
observability~\citep{hausknecht2015drqn}.  We use the
RecurrentPPO~\citep{schulman2017ppo} implementation from
\texttt{sb3-contrib}~\citep{sb3} for all three treatments, holding the
recurrent-layer architecture fixed and varying only the upstream
encoder.

\paragraph{Statistical reporting in deep RL.}
\citet{agarwal2021rliable} catalogued the field's practice of reporting
single-seed learning curves and showed how strongly that practice
overstates effect sizes.  Their \texttt{rliable} library has become
standard for reporting interquartile means with stratified-bootstrap
confidence intervals, and we adopt it throughout.  Our pre-registered
plan commits to IQM as the primary point estimate and 95\% bootstrap
CIs with 2,000 resamples as the primary uncertainty estimate.
